\documentclass[letterpaper]{article} %único necesario para: crear documento, hacer título, espacios, interlineado

\usepackage{vmargin} %necesario para ajustar las márgenes con \setmargins
%\usepackage{euler} %usa la fuente "euler" para las ecuaciones
\usepackage{amsfonts} %para Z estilizada
\usepackage{pifont} %para más estilos de viñetas
\usepackage{amsmath} %necesario equation*
\usepackage{amsthm} %necesario para proof enviroment
\usepackage{dsfont} %alguna letra caligráfica
\usepackage{tipa} %Epsilon bonito
\usepackage{graphicx} %para manejar imágenes
\usepackage{wrapfig} %para imágenes en modo "wrap"
\usepackage{lipsum} %para usar el texto de relleno
\usepackage{bold-extra} %permite usar textsc
\usepackage{tikz-cd} %para los diagramas conmutativos
\usepackage{mathrsfs} %para estilo especial de letra

\counterwithin{equation}{subsection}
\newtheorem{Tma}{Teorema}
\newtheorem{Prop}{Proposición}
\newtheorem{Def}{Definición}
\newtheorem{Lema}{Lema}
\newtheorem{Cor}{Corolario}
\newtheorem{Ejc}{Ejercicio}
\newtheorem{Ejm}{Ejemplo}
\newtheorem{Not}{Notación}

\newcommand{\Sh}{\normalfont{\text{Sh}}}
\newcommand{\PreSh}{\normalfont{\text{PreSh}}}
\newcommand{\PSu}{\normalfont{\text{P-Su}}}
\newcommand{\QSu}{\normalfont{\text{Q-Su}}}
\newcommand{\RSu}{\normalfont{\text{R-Su}}}
\newcommand{\Pgerm}{\normalfont{\text{P-germ}}}
\newcommand{\Qgerm}{\normalfont{\text{Q-germ}}}
\newcommand{\Rgerm}{\normalfont{\text{R-germ}}}
\newcommand{\germ}{\normalfont{\text{germ}}}
\newcommand{\Su}{\normalfont{\text{Su}}}
\newcommand{\Set}{\normalfont{\textbf{Set}}}
\newcommand{\Top}{\normalfont{\textbf{Top}}}
\newcommand{\Bund}{\normalfont{\textbf{Bund}}}
\newcommand{\Etale}{\normalfont{\textbf{Etale}}}
\newcommand{\Open}{\mathcal{O}}
\newcommand{\subseteqab}{\stackrel{ab}{\subseteq}}

\renewcommand*{\proofname}{\textbf{Prueba}} %Muestra "prueba" en lugar de "proof"
\renewcommand\refname{Referencias} %Muestra "Referencias" en lugar de "References"

%letterpaper: 21.59cm x 27.94cm
\setmargins{2.8 cm} %margen izquierdo
{1.5 cm} %margen Superior
{15.5cm} %anchura del texto
{23cm} %altura del texto
{0pt} %altura de los encabezados
{1.5cm} %espacio entre el texto y los encabezados
{5 cm} %altura del pie de página
{1 cm} %espacio entre el texto y el pie de página

\addtolength{\skip\footins}{2pc}

%\pagenumbering{gobble} %Quita la numeración de las páginas
\renewcommand{\baselinestretch}{1.5} %Aumenta el interlineado a n veces el automático
\relpenalty=9999 %Evita que se rompan las ecuaciones en el cambio de renglón
\binoppenalty=9999
\renewenvironment{abstract}
{\small
   \list{}{%
      \setlength{\leftmargin}{1.4cm}% <---------- CHANGE HERE
      \setlength{\rightmargin}{\leftmargin}%
   }%
   \item\relax}
{\endlist}

\renewcommand{\abstractname}{\textsc{Resumen}} %Muestra "Resumen" en lugar de "Abstract"

\usepackage{hyperref}
\hypersetup{
    colorlinks=true,
    linkcolor=blue,
    allcolors=blue,
    bookmarks=true,
    bookmarksopen=true,
    bookmarksnumbered=true
}
\usepackage{bookmark}

\date{}

\begin{document}

   \normalsize{
    \noindent\textit{Estudio sobre haces (2026)}

    \vspace{-0.1cm}
    \noindent Viernes 20 de febrero
    
    \vspace{-0.2cm}
}
\noindent\rule{10cm}{0.1pt}\\
\vspace{0.5cm}
\begin{center}
   \textsc{
      \Large{\textbf{Varias definiciones funtoriales de haz}}
   }\\ \vspace{0.7cm}
   %\large{\textsc{Juan Camilo Lozano Suárez \footnote{Estudiante de pregrado en matemáticas, Universidad Nacional de Colombia.\\Email: jclozanos@unal.edu.co}}}\\ \vspace{0.5cm}
\end{center}
\hspace{3.5cm} \hrulefill \hspace{3.5cm}

   \vspace{1cm}

   %\begin{abstract}
    %  \input{abstract}
     %   
      %\vspace{0.3cm}
      %\textsc{Palabras Clave.} Haz; espacio étalé; prehaz; homeomorfismo local; manojo; hacificación; equivalencia de categorías; local vs global.
   %\end{abstract}

   La Definición \ref{def:Godement} es debida a \cite[p.~109]{TATF}. 

   \begin{Def}\label{def:Godement}
      Sean $X$ un espacio topológico y $\mathscr{F}$ un prehaz de conjuntos sobre $X$. Se dice que $\mathscr{F}$ es un ${\text{haz}_1}$ de conjuntos cuando se cumplen las siguientes condiciones\footnote{Dados abiertos $U$ y $V\subseteq U$, y un $s\in\mathscr{F}(U)$, la imagen de $s$ por la aplicación $\mathscr{F}(U)\to\mathscr{F}(V)$ es la restricción de $s$ a $V$.}:
      \begin{itemize}
         \item[(F1):] Sea $(U_i)_{i\in I}$ una familia de conjuntos abiertos de $X$, $U$ la unión de los $U_i$ y $s',s''$ dos elementos de $\mathscr{F}(U)$; si las restricciones de $s'$ y $s''$ a cada $U_i$ son iguales, entonces $s'=s''$.
         \item[(F2):] Sea $(U_i)_{i\in I}$ una familia de conjuntos abiertos de $X$, con unión $U$, y supóngase dados unos $s_i\in \mathscr{F}(U_i)$ de tal manera que, cualesquiera que sean $i,j\in I$, las restricciones de $s_i$ y $s_j$ a $U_i\cap U_j$ sean iguales; entonces existe un $s\in \mathscr{F}(U)$ cuya restricción a $U_i$ es $s_i$ para todo $i\in I$.
      \end{itemize}
   \end{Def}

   \cite[p,~14]{ST} llama ``monoprehaces" o ``haces separados" a los prehaces que satisfacen \textit{(F1)}, y llama a \textit{(F2)} la ``condición de pegado". Nótese que por \textit{(F1)} se tiene que el elemento $s\in \mathscr{F}(U)$ cuya existencia garantiza \textit{(F2)} es único. Una definición alternativa de haz elimina la condición \textit{(F1)} pero exige la unicidad del $s\in \mathscr{F}(U)$ en \textit{(F2)}:

   \begin{Def}
      Sean $X$ un espacio topológico y $\mathscr{F}$ un prehaz de conjuntos sobre $X$. Se dice que $\mathscr{F}$ es una $\text{haz}_2$ de conjuntos si para cualesquiera $U\subseteqab X$, $(U_i)_{i\in I}$ cubrimiento abierto de $U$ y para cualquier colección $(s_i)_{i\in I}$ con $s_i\in \mathscr{F}(U_i)$ para cada $i\in I$, tal que\footnote{Denotamos $f^i_{ij}$ a la aplicación $\mathscr{F}(U_i)\to\mathscr{F}(U_i\cap U_j)$ y $f^j_{ij}$ a la aplicación $\mathscr{F}(U_j)\to\mathscr{F}(U_i\cap U_j)$.} $f^i_{ij}(s_i)=f^j_{ij}(s_j)$ para cualesquiera $i,j\in I$, existe un único $s\in \mathscr{F}(U)$ tal que\footnote{Denotamos $f_i$ a la aplicación $\mathscr{F}(U)\to\mathscr{F}(U_i)$.} $f_i(s)=s_i$ para todo $i\in I$. 
   \end{Def}

   La Definición \ref{def:Moerdijk} es debida a \cite[p.~66]{SGL}; para ella introducimos antes tres funciones canónicas.

   \begin{Def} 
      Sea $P$ un prehaz sobre un espacio topológico $X$. Dados $U\in\mathcal{O}(X)$ y $\left\lbrace U_i\right\rbrace_{i\in I}$ un cubrimiento abierto de $U$, definimos las siguientes funciones\footnote{Dados abiertos $U$ y $V\subseteq U$, y un $f\in PU$, representamos por $f|^U_V$ a la imagen de $f$ por la aplicación $PU\to PV$.}:
      \begin{itemize}
         \item $e:PU\to \prod_{i\in I}PU_i$ que a cada $f\in PU$ le asigna la familia $\left\lbrace f|^{U}_{U_i}\right\rbrace_{i\in I}$.
         \item $\pi_1:\prod_{i\in I}PU_i\to \prod_{(i,j)\in I\times I}P(U_i\cap U_j)$ que a cada $\left\lbrace f_i\right\rbrace_{i\in I}\in\prod_{i\in I}PU_i$ le asigna la familia $\left\lbrace {f_i}|^{U_i}_{U_i\cap U_j}\right\rbrace_{(i,j)\in I\times I}$.
         \item $\pi_2:\prod_{i\in I}PU_i\to \prod_{(i,j)\in I\times I}P(U_i\cap U_j)$ que a cada $\left\lbrace f_i\right\rbrace_{i\in I}\in\prod_{i\in I}PU_i$ le asigna la familia $\left\lbrace {f_j}|^{U_j}_{U_i\cap U_j}\right\rbrace_{(i,j)\in I\times I}$.
      \end{itemize}
      Llamamos a $e$ la función canónica de $PU$ en $\prod_{i\in I}PU_i$ y a $\pi_1$ y $\pi_2$ las funciones canónicas de $\prod_{i\in I}PU_i$ en $\prod_{(i,j)\in I\times I}P(U_i\cap U_j)$.
   \end{Def}

   \begin{Def}\label{def:Moerdijk}
      Un ${haz}_3$ de conjuntos $P$ sobre un espacio topológico $X$ es un prehaz de conjuntos tal que para cualquier $U\in \mathcal{O}(X)$ y cualquier cubrimiento abierto $\left\lbrace U_i\right\rbrace_{i\in I}$ de $U$, el siguiente es un diagrama igualador:
      \input{Diag1.tex}
   donde $e,\pi_1$ y $\pi_2$ son las respectivas funciones canónicas.
   \end{Def}

   \newpage
   \nocite{*}
   \bibliographystyle{apalike}
   \bibliography{refs.bib}
\end{document}
